% Generated draft: requires mathematical and source review.
\documentclass[11pt]{article}
\usepackage[margin=25mm]{geometry}
\usepackage{fontspec}
\setmainfont{Latin Modern Roman}
\usepackage{amsmath,amssymb,mathtools,unicode-math}
\setmathfont{Latin Modern Math}
\usepackage{enumitem,needspace,xurl,hyperref}
\setlength{\parindent}{0pt}
\setlength{\parskip}{5pt}
\setlength{\emergencystretch}{3em}
\newcommand{\N}{\mathbb{N}}
\newcommand{\Z}{\mathbb{Z}}
\newcommand{\Q}{\mathbb{Q}}
\newcommand{\R}{\mathbb{R}}
\newcommand{\C}{\mathbb{C}}
\newcommand{\F}{\mathbb{F}}
\newcommand{\A}{\mathbb{A}}
\newcommand{\PP}{\mathbb{P}}
\newcommand{\E}{\mathbb{E}}
\newcommand{\eps}{\varepsilon}
\newcommand{\dd}{\,\mathrm{d}}
\begin{document}
\section*{UnsolvedMath ID: 700688}
\textbf{Source problem code:} catalogue entry
\subsection{Definitions and mathematical statement}
A formation of finite groups is a nonempty class \(\mathcal F\) of finite groups, closed under isomorphisms, homomorphic images, and finite subdirect products. Here a subdirect product of \(G_1,\ldots,G_t\), with \(t\geq 1\), is a subgroup of \(\prod_{i=1}^t G_i\) whose projection onto each factor is surjective. A subformation of \(\mathcal F\) is a subclass \(\mathcal H\subseteq\mathcal F\) that is itself a formation; equality is permitted.

The source asks: enumerate all formations of finite groups all of whose subformations are \(S_n\)-closed. The supplied record does not define \(S_n\)-closed. Consequently, the following explicit formulation is conditional on interpreting this notation as closure under taking subnormal subgroups; this interpretation requires source verification.

A subgroup \(K\leq G\) is subnormal in \(G\) if there are an integer \(r\geq 0\) and subgroups
\[
K=K_0\trianglelefteq K_1\trianglelefteq\cdots\trianglelefteq K_r=G.
\]
Under the stated interpretation, the problem is to classify all formations \(\mathcal F\) of finite groups such that
\[
\begin{gathered}
\text{for every formation }\mathcal H\subseteq\mathcal F,
\text{ every }G\in\mathcal H,\text{ and every }K\leq G,\\
K\text{ subnormal in }G\quad\Longrightarrow\quad K\in\mathcal H.
\end{gathered}
\]
If the source uses a different meaning of \(S_n\)-closed, this conditional formulation must be revised. In particular, the supplied record does not specify an integer parameter \(n\) or a quantifier over such a parameter.
\subsection{Short English statement}
Classify the formations of finite groups for which every subformation has the source's \(S_n\)-closure property. If this means closure under subnormal subgroups, every subnormal subgroup of every group in each subformation must belong to that same subformation. The meaning of \(S_n\)-closed remains to be verified.
\subsection{Sources}
\begin{itemize}
\item Supplied local UnsolvedMath registry record, permanent ID \(700688\), exact source code \texttt{catalogue entry}, titled \emph{Kourovka Notebook Problem 10.73}. The record attributes the problem to L. A. Shemetkov, proposed in 1986, and labels it \emph{open}; that label is supplied registry data, not an independently verified mathematical assessment. The record was updated on \texttt{2026-10-16T14:00:00Z}.
\item \emph{The Kourovka Notebook (Unsolved Problems in Group Theory)}, original problem 10.73, 10th issue (1986); the supplied background identifies the 21st issue edition (2026), \texttt{arXiv:1401.0300v48}. The collection description names E. I. Khukhro and V. D. Mazurov as editors and describes coverage of Issues 1--21 (1965--2026). This attribution and edition information are supplied by the record; the cited source has not been independently consulted.
\end{itemize}
\end{document}
